\chapter{Requirements Analysis}
\label{ch:requirements}

\section{Functional Requirements}
The functional requirements below are derived directly from the behaviour
implemented in the source code.

\begin{enumerate}[nosep]
    \item The system shall allow a new user to register with a username
    (3--50 characters), an email address, and a password of at least 6
    characters, confirmed twice.
    \item The system shall reject registration when the supplied email is
    not a valid email address, or when the username or email is already in
    use.
    \item The system shall hash passwords with \texttt{password\_hash()}
    before storing them, and never store plaintext passwords.
    \item The system shall allow a registered user to log in with their
    email and password, verified with \texttt{password\_verify()}, and shall
    regenerate the session ID on successful login.
    \item The system shall allow a logged-in user to log out, which clears
    and destroys the server-side session.
    \item The system shall allow a logged-in user to start a game at one of
    two difficulties: Beginner ($9\times9$, 10 mines) or Advanced
    ($16\times16$, 40 mines).
    \item The system shall guarantee that the first cell revealed by the
    player, and all of its immediate neighbours, never contain a mine.
    \item The system shall reveal a cascade of adjacent zero-value cells
    automatically when a zero-adjacency cell is revealed (flood reveal).
    \item The system shall allow the player to flag and unflag unrevealed
    cells, up to the total mine count for the active difficulty.
    \item The system shall allow the player to chord on a revealed numbered
    cell when the number of adjacent flags equals its adjacent mine count,
    automatically revealing the remaining unflagged neighbours.
    \item The system shall allow the player up to 3 hints per game, each
    highlighting a random unrevealed, non-mine cell for 3 seconds.
    \item The system shall grant the player exactly one shield per game,
    which, once activated, prevents the next mine reveal from ending the
    game.
    \item The system shall end the game when the player reveals a mine
    (without an active shield) or when the 120-second timer reaches zero,
    and shall reveal all mines at that point.
    \item The system shall automatically submit the completed game's
    metrics (difficulty, time taken, cells revealed, hints used) to the
    backend once, and shall not re-submit the same completed game.
    \item The system shall recompute the score on the server from the
    submitted metrics rather than accept a client-supplied score, and shall
    reject metric values that fall outside the valid range for the given
    difficulty.
    \item The system shall reject a duplicate submission of the same result
    signature (user, difficulty, time, cells, hints) received within 10
    seconds of the previous one.
    \item The system shall display, for a logged-in user: personal
    statistics (games played, high score, average score, longest survival,
    most cells revealed), a history of up to 200 past games, and a global
    top-10 leaderboard, each filterable by difficulty.
\end{enumerate}

\section{Non-Functional Requirements}
\begin{itemize}[nosep]
    \item \textbf{Security} -- passwords must never be stored in plaintext;
    all SQL queries that include user input must use PDO prepared
    statements; session cookies must be HTTP-only; cross-origin requests
    must be restricted to an explicit allow-list of origins.
    \item \textbf{Data Integrity} -- the \texttt{users} table must guarantee
    unique usernames and unique email addresses; every stored game result
    must be linked to a valid user via a foreign key.
    \item \textbf{Usability} -- the interface should allow a player to start
    a game, understand their remaining time, mines, flags, and score, and
    review a completed game's outcome, without external instructions.
    \item \textbf{Maintainability} -- difficulty parameters (grid size, mine
    count) should be defined in a single, centralized configuration object
    on the frontend rather than duplicated across components.
\end{itemize}

\section{User Requirements}
A player needs an account, a way to start and play a timed game at a chosen
difficulty, visible feedback on their current score and remaining time and
mines, a way to review their past results, and a way to compare their best
scores against other players.

\section{Requirements Traceability}
Table~\ref{tab:req-trace} maps a representative subset of the functional
requirements above to the test cases used to validate them in
Chapter~\ref{ch:testing}.

\begin{table}[h]
\centering
\caption{Representative requirements-to-test-case mapping.}
\label{tab:req-trace}
\begin{tabular}{@{}p{2.2cm}p{9cm}p{1.5cm}@{}}
\toprule
\textbf{Req.\ \#} & \textbf{Requirement} & \textbf{Test ID} \\
\midrule
FR-3  & Password hashing on registration & T-2 \\
FR-7  & First-click safety & T-4 \\
FR-8  & Flood-fill cascade & T-5 \\
FR-10 & Chording precondition & T-6 \\
FR-12 & Shield blocks next mine reveal & T-7 \\
FR-15 & Server-side score recomputation & T-9 \\
FR-16 & Duplicate-submission rejection & T-10 \\
\bottomrule
\end{tabular}
\end{table}
